% ============================================================
%  roadmap_slides.tex
%  Supervisor meeting: the route from the manuscript's unresolved
%  Experiment 2 to the current state.
% ============================================================
\documentclass[11pt,aspectratio=169]{beamer}

\usetheme{Boadilla}
\usecolortheme{seahorse}
\setbeamertemplate{navigation symbols}{}
\setbeamertemplate{footline}[frame number]
\setbeamerfont{frametitle}{size=\large}
\setbeamercolor{block title}{bg=structure!20}

\usepackage{amsmath,amssymb}
\usepackage{booktabs}
\usepackage{array}

\newcommand{\good}[1]{\textbf{#1}}
\newcommand{\Prob}{\mathbb{P}}

\title[Queueing with RL: the route]{Multi-server queueing with reinforcement learning}
\subtitle{The route from the manuscript's Experiment~2 to the current state}
\author{}
\date{}

\begin{document}

\begin{frame}
\titlepage
\end{frame}

% ------------------------------------------------------------
\begin{frame}{Where the manuscript stops}

Table~3 of the manuscript, Experiment~2 (two flexible servers, $3\times$ cost
asymmetry):

\begin{center}
\begin{tabular}{lcc}
\toprule
Policy & $\bar g \cdot \Lambda$ & ratio to RVI \\
\midrule
RVI (optimal)      & 40.79  & 1.00 \\
FIFO               & 70.38  & 1.73 \\
\midrule
DCL from FIFO      & 180.3  & 4.42 \\
DCL from StochFIFO & 325.7--380.8 & 7.99--9.34 \\
\bottomrule
\end{tabular}
\end{center}

\vfill
\begin{block}{The manuscript's own note}
\emph{``DCL however fails to capture to learn the desired behaviour.
To be fixed in the future.''}
\end{block}

\vfill
That is the starting point.  The finding since: this was not a tuning problem.

\end{frame}

% ------------------------------------------------------------
\begin{frame}{The route}

\begin{center}
\begin{tabular}{cp{9.2cm}}
\toprule
\textbf{1} & Exhaust DCL --- diagnose why it has a fixed point \\[2pt]
\textbf{2} & Switch to PPO $\to$ \good{five failure mechanisms}, each isolated \\[2pt]
\textbf{3} & Queue lateness as an alternative objective $\to$ the failure \good{inverts} \\[2pt]
\textbf{4} & One diagnosis connects most mechanisms $\to$ \good{reformulate the action space} \\[2pt]
\textbf{5} & Fraction served late $\to$ ill-posed $\to$ \good{escalation rule} \\[2pt]
\textbf{6} & Scaling $\to$ turns out to be an \good{instance-design} problem \\
\bottomrule
\end{tabular}
\end{center}

\vfill
Throughout: RVI as the exact benchmark, and every claim reported as a
\emph{seed distribution} (``$k$ of 8 seeds within $x\%$'') rather than a
single run.

\end{frame}

% ------------------------------------------------------------
\begin{frame}{1. Why DCL gets stuck}

DCL trains a \emph{classifier} on best actions identified by rollouts.  No value
function, no discounting, undiscounted horizon $H$.

\vfill
\begin{block}{The fixed point}
Started from a deterministic FIFO base, the states visited are FIFO-reachable
states.  On \emph{FIFO's own} distribution the deviations do not look
advantageous, so generation 1 reproduces FIFO --- and generation 2 starts from
the same policy.
\end{block}

\vfill
The manuscript already mitigates this (stochastic FIFO base, $\varepsilon$-greedy
wrapper).  It was not enough on Experiment~2.

\vfill
\emph{No convergence guarantees exist here:} DCL is approximate policy iteration
with classification error, and the standard guarantees need that error
controlled uniformly.

\end{frame}

% ------------------------------------------------------------
\begin{frame}{2. PPO: five failure mechanisms}

Each fix removed from the otherwise complete recipe, 8 seeds, Experiment~2.

\begin{center}
\small
\begin{tabular}{p{3.4cm}p{5.4cm}cc}
\toprule
\textbf{Mechanism} & \textbf{What goes wrong} & \textbf{Without} & \textbf{Failure} \\
\midrule
Semi-MDP discounting & decisions separated by a variable number of periods,
including zero; discounting each as if a period passed subsidises idling
& \good{0/8} & never-serve \\[2pt]
Value-target scale & returns $O(10^2$--$10^3)$ vs advantages $O(1)$; value loss
monopolises the shared trunk & \good{0/8} & bimodal \\[2pt]
Data distribution & one deep-backlog excursion moves all environments out of
the deployment region & \good{1/8} & never-serve \\[2pt]
Reward sparsity & binary cost flat before the deadline & \good{2/8} & never-serve \\[2pt]
Extraction & deployed policy is the \emph{argmax}; training certifies the
\emph{stochastic} policy & \good{3/8} & argmax $\gg$ stoch \\
\midrule
\textbf{Full recipe} & & \good{6/8} & \\
\bottomrule
\end{tabular}
\end{center}

\vfill
\small
Both attractors recur \emph{exactly} in every failed cell: never-serve at
$9.876\times$, FIFO clone at $1.738\times$ the optimum.

\end{frame}

% ------------------------------------------------------------
\begin{frame}{A note on the reward: the objective did not change}

\begin{columns}[T]
\column{0.48\textwidth}
\textbf{Evaluated} \\[2pt]
The manuscript's binary FIL lateness
\[ \tfrac{c_n}{\gamma}\mathbf{1}\{W_n > D\} \]
RVI benchmark unchanged.

\column{0.48\textwidth}
\textbf{Trained on} \\[2pt]
Same cost $+$ a potential
\[ \Phi_n(t) = c_n \tfrac{\min(t,D)}{D} \]
accrued per pre-deadline tick, refunded on service.
\end{columns}

\vfill
\begin{block}{Why the benchmark survives}
$r' = r + \Phi(s) - \Phi(s')$ telescopes, so every stationary policy keeps its
long-run average.  Verified numerically: FIFO and RVI evaluate identically with
and without shaping, to four decimals.
\end{block}

\vfill
Effect: dense gradient before the deadline.  $2/8 \to 6/8$ seeds.

\end{frame}

% ------------------------------------------------------------
\begin{frame}{3. Queue lateness: the failure inverts}

Motivated by the question in your comments --- would queue lateness be better for
shaping?

\vfill
\begin{columns}[T]
\column{0.48\textwidth}
\textbf{What gets better} \\[2pt]
The extraction gap disappears: argmax $\approx$ stochastic everywhere.  Cleaner
threshold structure.

\column{0.48\textwidth}
\textbf{What gets worse} \\[2pt]
Per-tick cost is $\Theta(B^2)$ in the excess $B$.  Rollout reward rates swing
$-380$ to $-7500$ \emph{within healthy runs}.
\end{columns}

\vfill
\begin{block}{Consequence}
The tuned recipe converges to a \emph{bit-exact FIFO clone on 8 of 8 seeds}.
A different configuration is required: average-reward (differential) objective
with a $64\times$ larger batch.
\end{block}

\vfill
\small
It also does \emph{not} fix sparsity: queue lateness is equally flat before the
deadline.  And the cost is no longer tick-rate invariant --- FIFO measures
$1781 / 3035 / 5266$ at $\gamma = 1.5/3/6$ (binary: $74/71/70$).

\end{frame}

% ------------------------------------------------------------
\begin{frame}{Why the annealing guard breaks under heavy tails}

The guard lets the behaviour temperature $T$ fall only while an EMA of the
reward rate does not degrade, measured against a \emph{ratchet} reference (best
EMA so far).

\vfill
\begin{columns}[T]
\column{0.47\textwidth}
\textbf{Why a ratchet?} \\[2pt]
Without it a slowly degrading policy drags the reference down with it --- each
small decline stays within tolerance --- and the guard keeps sharpening a
collapsing policy.

\column{0.47\textwidth}
\textbf{Its flaw under QL} \\[2pt]
One lucky quiet stretch sets an unrepeatable reference.  It is stored forever,
so the normal level then reads as ``degrading''.
\end{columns}

\vfill
Result: \good{4 of 8 seeds finish 1000 updates still at $T = 1.0000$} --- and
those are exactly the bad seeds.

\vfill
\begin{block}{Leaky ratchet}
On each \emph{failed} check, relax the reference $5\%$ toward the current EMA.
Separation of speeds: temperature backs off \emph{immediately} on degradation,
the reference forgets only \emph{slowly} --- so genuine decline is still caught.
Engagement $4/8 \to 7/8$.
\end{block}

\end{frame}

% ------------------------------------------------------------
\begin{frame}{4. The diagnosis that connects them}

Why is the argmax the problem?  A skipped job \emph{returns at the next tick}.

\vfill
With serve-probability $p$ and $m$ further presentations before the deadline:
\[
  \Prob(\text{served in time}) \;=\; 1-(1-p)^{m+1}
  \qquad\text{e.g. } p = 0.4,\ m = 8 \;\Rightarrow\; 0.99
\]

\vfill
\begin{block}{Re-presentation forgiveness}
The \emph{stochastic} policy is behaviourally near-optimal even where the logit
sign is wrong.  The advantage vanishes, so training never corrects it --- but the
\emph{argmax} turns exactly those states into deterministic errors, at every
re-presentation.
\end{block}

\vfill
Most of the stabilisation apparatus --- temperature annealing, snapshots, readout
selection, seed selection --- existed only to contain this one mechanism.

\end{frame}

% ------------------------------------------------------------
\begin{frame}{The per-event action space}

\begin{columns}[T]
\column{0.47\textwidth}
\textbf{Sequential (manuscript)} \\[2pt]
Ordered candidate list; binary accept/skip at each position; a skipped job
returns next tick.

\column{0.47\textwidth}
\textbf{Per-event} \\[2pt]
Each idle capacity unit makes \emph{one} $(N{+}1)$-way decision: idle, or serve
type $a$.  Infeasible types masked.
\end{columns}

\vfill
\begin{block}{This is what your comment on \S3.2 proposed}
\emph{``kan een event tot meer dan 1 service start leiden? indien max 1 dan kies
je de beste uit de toegelaten toewijzingen''}
\end{block}

\vfill
Same physical policies --- verified by three equivalence gates:

\begin{center}
\small
\begin{tabular}{lcc}
\toprule
 & sequential & per-event \\
\midrule
FIFO, Exp 2         & 70.3802  & 70.3802 \quad(bit-exact) \\
never-serve         & 399.9782 & 399.9782 \quad(bit-exact) \\
RVI optimum         & 40.786   & 40.785 \quad(solver tolerance) \\
\bottomrule
\end{tabular}
\end{center}

\vfill
\small Only the \emph{credit structure} changes.  Per-decision action set is
$O(N)$ regardless of the number of servers.

\end{frame}

% ------------------------------------------------------------
\begin{frame}{The apparatus becomes unnecessary --- and sometimes harmful}

Experiment 2, binary cost, seeds within 5\% of RVI:

\begin{center}
\begin{tabular}{lcc}
\toprule
 & sequential & per-event \\
\midrule
with temperature annealing    & 6/8 & 4/8 \\
\good{without} annealing      & 3/8 & \good{8/8, exact optimum} \\
\bottomrule
\end{tabular}
\end{center}

\vfill
\begin{block}{But it is problem-dependent}
On Experiment~3 the dissociation runs the other way: 7/8 \emph{with} annealing,
3/8 without.
\end{block}

\vfill
The honest statement: a stabiliser is matched to a failure mode, and outlives
its usefulness with it.  It should be a switchable ingredient, not a fixed part
of the recipe.

\end{frame}

% ------------------------------------------------------------
\begin{frame}{5. The objective you actually want: fraction served late}

Charging at departure gives delayed credit \emph{and} creates the incentive the
manuscript already identified (\S3): once late, serving earns nothing.

\vfill
Instead charge $c_n$ \emph{at the deadline crossing} --- same total on every
sample path, at the decision-relevant moment.  In the FIL projection:
\[
  r_4 \;=\; c_n\Big[\underbrace{\mathbf{1}\{W = D+1\}}_{\text{the FIL crosses}}
  \;+\; \underbrace{\tfrac{\lambda_n}{\gamma}\,\mathbf{1}\{W \ge D+2\}}_{\text{expected crossings behind it}}\Big]
\]

\vfill
\begin{block}{Why $\lambda_n/\gamma$ and not a probability}
One \emph{slot} crosses per tick, and it may hold $0,1,2,\dots$ jobs.  We count,
so we need the \emph{mean}.  Check: if the queue stays backed up every arrival
eventually crosses, so crossings per unit time must equal $\lambda_n$ ---
and $\gamma \cdot (\lambda_n/\gamma) = \lambda_n$ for every $\gamma$.
\end{block}

\end{frame}

% ------------------------------------------------------------
\begin{frame}{The escalation rule}

\begin{alertblock}{The objective is ill-posed on the FIL projection}
Letting a job rot \emph{suppresses that type's arrival stream} --- arrivals are
only tracked when the queue is empty.  So rotting is not merely free, it is
profitable.  RVI discovers this: the computed optimum drifts below the
achievable cost, and learned policies ``beat'' the benchmark.
\end{alertblock}

\vfill
\textbf{Remedy: a modelling rule, not a reward adjustment.}  When a capacity
unit can serve a job already past its deadline, the assignment is forced
(oldest-late first, non-preemptive) --- an SLA escalation rule.

\vfill
Implemented purely by restricting the admissible action set to a singleton, so
PPO, DCL and RVI all solve the restricted MDP \emph{with no algorithmic change}.

\vfill
\begin{columns}[T]
\column{0.47\textwidth}
Degeneracy removed by construction.
\column{0.47\textwidth}
Late states become transient $\Rightarrow$ RVI truncation converges
(bit-identical at $M = 30, 60, 90$).
\end{columns}

\end{frame}

% ------------------------------------------------------------
\begin{frame}{Three different ``deep queue'' problems}

Easily conflated --- they have different causes and different remedies.

\vfill
\begin{center}
\small
\begin{tabular}{p{3.2cm}p{4.6cm}p{3.0cm}p{2.6cm}}
\toprule
\textbf{Problem} & \textbf{What happens} & \textbf{Remedy} & \textbf{Applies to} \\
\midrule
Data poisoning & deep excursions dominate the training data; the deployment
region disappears from it & staggered environment resets & all objectives \\[3pt]
Return variance & deep queues make returns heavy-tailed & average-reward $+$
large batch & queue lateness \\[3pt]
Abandonment & letting a job rot is profitable & escalation rule & \good{tardiness only} \\
\bottomrule
\end{tabular}
\end{center}

\vfill
\begin{block}{Why escalation is not needed for lateness}
Under binary and queue lateness the cost \emph{keeps accruing} while a job is
late --- quadratically under QL.  Clearing late work is rewarded by the objective
itself.  \emph{All queue-lateness results (8/8 on both instances) were obtained
without escalation.}
\end{block}

\end{frame}

% ------------------------------------------------------------
\begin{frame}{Where that leaves us: three objectives, two instances}

Seeds within 5\% of RVI, per-event action space, no per-cell tuning.

\begin{center}
\begin{tabular}{p{3.2cm}p{4.2cm}cc}
\toprule
\textbf{Objective} & \textbf{Needs} & \textbf{Exp 2} & \textbf{Exp 3} \\
\midrule
Binary / FIL lateness   & shaping; small batch      & 8/8 & 7/8 \\
Queue lateness          & average-reward; $64\times$ batch & 8/8 & 8/8 \\
Fraction late           & escalation rule           & 8/8 & 8/8 \\
\bottomrule
\end{tabular}
\end{center}

\vfill
\begin{block}{The recipe follows from the tail of the cost function}
Queue lateness is the only objective with unbounded per-tick cost
($\Theta(B^2)$), and correspondingly the only one needing the large-batch
average-reward configuration.
\end{block}

\vfill
\small
The benchmark itself needed correcting twice --- both times detected because a
\emph{learned policy evaluated below the computed optimum}, which can only mean
a loose benchmark.

\end{frame}

% ------------------------------------------------------------
\begin{frame}{6. Scaling turned out to be a modelling problem}

First attempt on a six-type chain: the learned policy beat the best heuristic by
only \good{2.9\%}.  Not under-training --- a $2.5\times$ budget gives the same
answer.

\vfill
\begin{alertblock}{Why}
On that instance cost, service rate and deadline all increase together, so the
$c\mu$ ordering \emph{coincides} with the urgency ordering.  The heuristic was
not clever; the instance was accommodating.
\end{alertblock}

\vfill
Screening four candidate structures at three types, where RVI is still
computable:

\begin{center}
\small
\begin{tabular}{lccc}
\toprule
Structure (3 types, load 2.0) & FIFO/RVI & $c\mu$/RVI & $\min$ \\
\midrule
\good{cost and deadline anti-aligned} & \good{1.49} & \good{1.77} & \good{1.49} \\
equal costs, deadlines heterogeneous  & 1.10 & 1.16 & 1.10 \\
non-interval skill overlap            & 1.10 & 1.36 & 1.10 \\
generalist bottleneck (rewards idling)& 1.21 & 1.01 & 1.01 \\
\midrule
existing chain (all aligned)          & 1.19 & 1.19 & 1.19 \\
\bottomrule
\end{tabular}
\end{center}

\end{frame}

% ------------------------------------------------------------
\begin{frame}{Scaling the winning structure}

Recipe unchanged, 8 seeds per instance.

\begin{center}
\begin{tabular}{lccccc}
\toprule
Instance & types & FIFO & $c\mu$ & PPO (8 seeds) & vs $c\mu$ \\
\midrule
anti-aligned & 3 & 1.372 & 1.626 & 0.929--0.965 & \good{$-43\%$} \\
anti-aligned & 4 & 1.518 & 1.837 & 1.124--1.201 & $-35$ to $-39\%$ \\
anti-aligned & 6 & 1.745 & 2.124 & 1.422--1.634 & $-23$ to $-31\%$ \\
\bottomrule
\end{tabular}
\end{center}

\vfill
At three types, where the optimum is computable ($\mathrm{RVI} = 0.920$):
\good{8/8 seeds within 1.0--4.9\%}.

\vfill
\begin{block}{The decisive comparison, at equal size}
Original six-type chain: $-2.9\%$ over the best heuristic. \\
Redesigned six-type instance: $-23$ to $-31\%$. \\[2pt]
Same algorithm, same recipe.  \emph{The earlier ceiling was a property of the
benchmark.}
\end{block}

\end{frame}

% ------------------------------------------------------------
\begin{frame}{Open questions}

\begin{enumerate}
  \item \textbf{A deadline-aware baseline is missing.}  FIFO and $c\mu$ are both
  deadline-blind, which makes them a weak comparison for a deadline objective.
  EDF / least slack, and a generalised $c\mu$ index, should be added.

  \vspace{4pt}
  \item \textbf{No exact reference beyond three job types.}  A fluid or LP lower
  bound would be needed to say how much is left on the table --- the margin
  narrows with size ($-43\%$, $-38\%$, $-31\%$) and we cannot currently tell
  whether that is the instance or the method.

  \vspace{4pt}
  \item \textbf{The network does not transfer across sizes.}  A fixed-width MLP
  assigns a separate semantic role to each input position.  Parameter sharing
  over job types and server pools (a bipartite encoder) is the natural next
  step, and is what a $10\times10$ claim would require.

  \vspace{4pt}
  \item \textbf{Compute is not the constraint.}  The six-type experiment cost
  $\approx$ 96 core-hours; the budget is 77{,}000.
\end{enumerate}

\end{frame}

\end{document}
